% Introduction. Claims: paper/CLAIMS.md; novelty: paper/NOVELTY.md
\section{Introduction}
\label{sec:intro}

After an earthquake, flood or fire, the ground robots sent to search a
disaster site need a mission-critical link for command and control (C2) and
video, and the terrestrial network they would use is often gone or
overloaded~\cite{murphy2014disaster}. A 5G base station (gNB) carried by an
unmanned aerial vehicle (UAV) can bring a private cell to the
robots~\cite{zeng2016wireless,mozaffari2019tutorial}. Such a cell has to be
managed: something must decide where to fly it, notice when a robot is about
to lose its link, and keep the links usable under the load the robots put on
them. On site there is no operator to tune it and no reliable map of terrain
that may have just changed, so the management has to be autonomous, proactive
and blind to the terrain, which is the setting of zero-touch
management~\cite{coronado2022zsm}.

Our earlier testbed study~\cite{makhubela2026satnac} showed that terrain decides
whether moving such a node helps: an uninformed roaming node raised a weak
robot's C2 availability in open field and cost about 22 points in forest. The
natural next step is a learning-based management layer that forecasts link
failures and places the UAV in three dimensions without being told the terrain.
This paper builds that layer as part of a three-loop management stack
(Fig.~\ref{fig:arch}), flies it in closed loop over real 5G software, and asks
what each loop delivers. The answer turned out to depend less on the learning
than on whether the managed system was measured and loaded soundly.

We test four claims, each against an explicit control (Section~\ref{sec:setup}):
\begin{itemize}
\item[\textbf{C1}] Long forest outages under robot video are caused by uplink
queuing inside the user equipment (UE), not by coverage, and a delay-adaptive
uplink cap at each robot removes them without starving strong links.
\item[\textbf{C2}] A terrain-blind model forecasts the onset of C2 failure
20\,s ahead on an unseen terrain better than chance, and higher-capacity models
do worse.
\item[\textbf{C3}] Terrain-blind models do not recover how terrain changes the
effect of altitude, so they cannot choose altitude.
\item[\textbf{C4}] Learned placement does not beat fixed hover, over repeated
missions in two terrains.
\end{itemize}
C1 holds; C2 and C4 hold weakly; C3 holds as a negative result.

\textit{Contributions.} The work falls under wireless and 5G networks,
mission-critical networks, proactive and autonomic management, artificial
intelligence (AI)-based management, monitoring and measurement, and testbed
experimentation. Beyond the testbed and terrain comparison
of~\cite{makhubela2026satnac} (Table~\ref{tab:satnac}), it contributes:
\begin{enumerate}
\item \emph{A diagnosis and an autonomic fix.} The forest outages are a silent
queueing failure of the managed system. A per-robot cap driven only by
round-trip time and bytes sent, with no change to the radio access network
(RAN), shortens outages from minutes to seconds and lifts both placement arms
flown with and without it (Section~\ref{sec:r-cap}).
\item \emph{A liveness monitor with a measurement-integrity check.} A second
silent failure, the UE process dying after radio-link failure, contaminated
12.4\,\% of training observations; we give a mechanistic detection rule and
validate it against the monitor's own record (Section~\ref{sec:dead}).
\item \emph{An evaluation of terrain-blind learning for proactive management}
on held-out terrain and in closed loop, with the negative results stated
(Sections~\ref{sec:r-forecast} and~\ref{sec:r-flip}).
\item \emph{A closed-loop account of learned placement}, showing why a single
mission is too short to learn the forest's altitude effect and why a cautious
learner reduces to hover (Section~\ref{sec:r-place}).
\end{enumerate}
All results come from emulation over real 5G software; a hardware trial is the
next step.
